%% ma: L10-B04-0002
%% lop: 10
%% chuong: 2
%% bai: 4
%% dang: 10.4.1
%% loai: TLN
%% muc: VD
%% dap_an: "1340"
%% nguon: "(NBV)-ĐỀ SỐ 12-MỖI NGÀY 1 ĐỀ THI 2025.pdf (D12, MỖI NGÀY 1 ĐỀ - Đề số 12), Phần III Câu 1 (Sở Vĩnh Phúc 2025)"
%% kien_thuc: "Bài 4 lớp 10 (hệ bất phương trình bậc nhất hai ẩn, miền nghiệm, tối ưu trên miền đa giác)"
%% trang_thai: chinh-thuc
%% ghi_chu: "Đáp án gốc 1340 đúng (kiểm bằng duyệt các cặp số nguyên: lãi lớn nhất 115500 nghìn đồng tại x = 250, y = 280). Đã bỏ ảnh minh họa xưởng may (không chứa dữ kiện toán). Cùng ý tưởng với dạng 10.4.1 (quy hoạch tuyến tính tối đa lợi nhuận), nhập vào dạng đã duyệt."
\begin{ex}
Một xí nghiệp may áo vest và quần âu để chuẩn bị cho dịp cuối năm. Biết may $1$ áo vest hết $2$ m vải và cần $20$ giờ; may $1$ quần âu hết $1{,}5$ m vải và cần $5$ giờ. Xí nghiệp được giao sử dụng không quá $920$ m vải và số giờ công không vượt quá $6400$ giờ. Theo khảo sát thị trường, số lượng quần âu bán ra không nhỏ hơn số lượng áo vest bán ra và số lượng quần âu bán ra không vượt quá $2$ lần số lượng áo vest bán ra. Khi xuất ra thị trường, $1$ chiếc áo vest lãi $350$ nghìn đồng, $1$ chiếc quần âu lãi $100$ nghìn đồng. Gọi $x$, $y$ lần lượt là số áo vest và quần âu xí nghiệp cần may và bán ra thị trường để xí nghiệp có số tiền lãi cao nhất. Tính giá trị của biểu thức $T=2x+3y$.
\shortans{1340}
\loigiai{Các ràng buộc: $2x+1{,}5y\le920$, $20x+5y\le6400$, $x\le y\le2x$, $x,y\ge0$, tức
\[\begin{cases}4x+3y\le1840,\\4x+y\le1280,\\x\le y\le2x.\end{cases}\]
Miền nghiệm là tứ giác $OABC$ với $O(0;0)$, $A(184;368)$ (giao của $y=2x$ và $4x+3y=1840$), $B(250;280)$ (giao của $4x+3y=1840$ và $4x+y=1280$), $C(256;256)$ (giao của $y=x$ và $4x+y=1280$).

Tiền lãi $L=350x+100y$ (nghìn đồng): $L(O)=0$, $L(A)=101200$, $L(B)=115500$, $L(C)=115200$. Lãi lớn nhất tại $B(250;280)$.

Vậy $T=2\cdot250+3\cdot280=1340$.}
\end{ex}
